\chapter{Другий вихід}
% full version: chapters/16_chemoutput.tex

Машина має два виходи. Швидкий --- м'язи. Повільний --- хімія тіла: мозок керує серцем, судинами, залозами й через кров розсилає повідомлення всьому тілу одразу. Кров --- найширокомовніша шина з усіх: одне повідомлення доходить до кожної клітини приблизно за хвилину, а зміст йому, як завжди, надає «замок» отримувача.

\section*{Газ і гальмо}

Серце б'ється саме, у нього свій метроном. Мозок лише підкручує темп двома проводами протилежного знака. \nt{NORA} --- педаль газу: розганяє, але повільно. \nt{ACET} --- гальмо: сповільнює швидко, бо йому не потрібен довгий ланцюжок хімічних реакцій. І тут нарешті знаходиться робота для типу \nt{ADRE}: адреналін, викинутий у кров, --- та сама педаль газу, тільки натиснута одразу для всього тіла.

\pencilfig{16}{\pencilart{16}}{Серце ведуть два проводи протилежного знака: газ і гальмо.}

\section*{Каскад зі зворотним зв'язком}

Багатьма гормонами керує каскад: мозок наказує першій залозі, та --- другій, друга виділяє гормон, а гормон повідомляє мозку, що його вже досить. Це термостат із трьох щаблів. Інженер знає, чим небезпечна така схема: якщо петлю зробити надто «чутливою», вона почне розгойдуватися --- рівень не триматиметься, а коливатиметься. Тому точність такого регулятора обмежена: зробити його водночас стійким і дуже точним неможливо.

\pencilfig{127}{%
% three identical first-order stages with proportional feedback: secretion = max(0, K (target - level)).
% K = 3 settles at 3/4 of the target; K = 12 (above the stability limit K = 8) keeps oscillating. x = 0.42 t, y = 0.55 level
\foreach \yo/\t in {1.75/{слабка петля: спокійно, але повз ціль}, 0/{сильна петля: ближче до цілі, але розгойдується}}{
  \draw[pen] (0,\yo) -- (8.5,\yo);
  \draw[pcGray, densely dashed, line width=0.5pt] (0,\yo+0.55) -- (8.5,\yo+0.55);
  \node[lbl, anchor=south west] at (2.0,\yo+0.8) {\t};}
\node[lbl, anchor=west] at (8.55,2.3) {ціль}; \node[lbl, anchor=west] at (8.55,0.55) {ціль};
\begin{scope}[yshift=1.75cm]
\draw[penlite=pcBlue] plot[smooth] coordinates {(0.00,0.00) (0.17,0.01) (0.34,0.08) (0.50,0.19) (0.67,0.33) (0.84,0.46) (1.01,0.55) (1.18,0.59) (1.34,0.58) (1.51,0.55) (1.68,0.49) (1.85,0.43) (2.02,0.37) (2.18,0.34) (2.35,0.33) (2.52,0.34) (2.69,0.37) (2.86,0.40) (3.02,0.43) (3.19,0.44) (3.36,0.45) (3.53,0.45) (3.70,0.44) (3.86,0.42) (4.03,0.41) (4.20,0.40) (4.37,0.39) (4.54,0.39) (4.70,0.40) (4.87,0.40) (5.04,0.41) (5.38,0.42) (5.88,0.42) (6.38,0.41) (7.06,0.41) (8.40,0.41)};
\end{scope}
\draw[penlite=pcRed] plot[smooth] coordinates {(0.00,0.00) (0.17,0.05) (0.34,0.30) (0.50,0.68) (0.67,1.02) (0.84,1.23) (1.01,1.30) (1.18,1.26) (1.34,1.16) (1.51,1.02) (1.68,0.87) (1.85,0.72) (2.02,0.58) (2.18,0.47) (2.35,0.39) (2.52,0.38) (2.69,0.46) (2.86,0.57) (3.02,0.67) (3.19,0.71) (3.36,0.70) (3.53,0.65) (3.70,0.58) (3.86,0.50) (4.03,0.43) (4.20,0.41) (4.37,0.45) (4.54,0.53) (4.70,0.61) (4.87,0.65) (5.04,0.64) (5.21,0.60) (5.38,0.54) (5.54,0.47) (5.71,0.42) (5.88,0.42) (6.05,0.48) (6.22,0.56) (6.38,0.62) (6.55,0.64) (6.72,0.62) (6.89,0.56) (7.06,0.50) (7.22,0.44) (7.39,0.42) (7.56,0.45) (7.73,0.53) (7.90,0.60) (8.06,0.63) (8.23,0.62) (8.40,0.58)};
\node[lbl, anchor=north] at (4.2,-0.05) {час};
}{Модель гормонального каскаду з трьох щаблів. Слабкий зворотний зв'язок заспокоює рівень, але лишає його помітно нижчим за ціль. Сильний підводить ближче, зате рівень уже не тримається, а гойдається.}

\section*{Імпульси, а не рівень}

Деякі гормони працюють лише імпульсами. Якщо подавати їх постійно, залоза-отримувач «утомлюється» й вимикається; якщо короткими імпульсами раз на годину --- працює. Отримувач читає не рівень, а ритм, --- як синапс із розділу про абетку Морзе, що чує лише зміни.

\section*{Доба}

І нарешті, у тілі є повільний генератор із періодом близько доби. Його природний період трохи відрізняється від двадцяти чотирьох годин, і щоранку світло підлаштовує його, як радіоприймач налаштовується на станцію. Не плутайте його з тактовим генератором комп'ютера: він не відмірює кроків обчислення. Він перемикає режими роботи всієї машини --- день і ніч, неспання й сон.

\fullversion{16}
